\bookchapter{第十章：移動支付與生活服務}

\chapterintro{說實話，我爸媽那一輩人對“手機付錢”這件事，一開始是抗拒的。“錢在手機裡，被人偷了怎麼辦？”這個擔心很正常，但現實是：你褲兜裡揣著的現金，丟的機率比手機支付被盜刷的機率高得多。手機支付不是把錢存進手機，它只是一把鑰匙，一把開啟你銀行賬戶的鑰匙。}

\booksection{一、到底怎麼繫結銀行卡？}

不管你是用支付寶還是微信，第一次綁卡都要走這幾步：

\begin{enumerate}
  \item 開啟App，找到"我的"或"錢包"。點"新增銀行卡"。
  \item 輸入卡號。現在大部分App都支援直接拍照識別卡號，不用一個個數字敲。
  \item 銀行會發一條驗證碼簡訊到你預留的手機號上。填進去，設定一個6位數字的支付密碼。
  \item 這個支付密碼和你的銀行卡密碼、鎖屏密碼，最好是三個不同的密碼，別圖省事設成一樣的。
\end{enumerate}

有個小竅門：如果擔心安全問題，可以去銀行App裡開一張II類或III類電子賬戶，專門用來繫結手機支付。這張卡裡只放平時零花的錢，就算出了事，損失也有限。

\booksection{二、付款碼是“誰看誰收錢”}

很多人搞不清楚付款碼和收款碼的區別，這個區別很重要。

\begin{itemize}
  \item \textbf{付款碼}：你展示給收銀員掃的那個條形碼+二維碼。這個碼相當於你的"電子錢包"，誰掃了誰就能從你賬上劃走錢。所以\textbf{千萬別截圖發給任何人}。這個碼每分鐘自動重新整理一次，就是為了防止被盜用。
  \item \textbf{收款碼}：別人掃了這個碼給你轉賬。這個碼可以截圖儲存，可以發給別人，安全。
\end{itemize}

認準這個區別，能避開很多詐騙陷阱。凡是有人讓你“把付款碼截圖發過來”的，百分之百是騙子。

\booksection{三、這些安全設定，花5分鐘設好能省很多心}

\begin{itemize}
  \item \textbf{指紋/面容支付}：在支付寶和微信的支付設定裡，把生物支付開啟。這樣每次付款要按指紋或掃臉，就算手機丟了，撿到的人也沒法花你的錢。
  \item \textbf{安全鎖}：支付寶和微信都有"安全鎖"功能。開啟之後，進入App的"錢包"頁面需要額外驗證一次身份。路徑：支付寶 我的 設定 賬號與安全 安全中心 安全鎖。
  \item \textbf{免密支付檢查}：很多人簽約了各種免密支付，自己都忘了。去支付寶和微信的"支付設定"裡，找到"免密支付/自動扣款"，看看簽約了哪些服務，把不需要的關掉。
  \item \textbf{登入裝置管理}：定期看看哪些裝置登入過你的賬號。如果發現不認識的裝置，立刻踢掉。路徑：支付寶 設定 賬號與安全 登入裝置管理。
\end{itemize}

\booksection{四、公交地鐵：手機就是交通卡}

現在出門坐公交地鐵，已經很少人用實體的交通卡了。手機刷閘機，有兩條路可以走：

\begin{itemize}
  \item \textbf{手機錢包裡的虛擬交通卡}：在蘋果錢包、華為錢包、小米錢包裡直接開通一張虛擬交通卡。這個方案的好處是：刷的時候不用解鎖手機，直接靠上去就行，響應速度最快，而且沒網也能刷。部分城市的虛擬交通卡還有和實體卡一樣的換乘優惠。
  \item \textbf{支付寶/微信的乘車碼}：在App裡搜尋"乘車碼"，開通你所在城市的公交或地鐵乘車碼。這個方案的好處是：不用充值，直接從餘額或繫結的銀行卡裡扣。但需要聯網，而且刷之前要開啟App找到碼。
\end{itemize}

我個人更推薦用手機錢包的虛擬交通卡。雖然第一次要充個10塊20塊進去，但用起來真的方便太多，尤其是趕地鐵的時候，不用停下來解鎖手機找App。

\booksection{五、生活繳費、外賣、打車}

\begin{itemize}
  \item \textbf{水電氣繳費}：在支付寶或微信的"生活繳費"裡，把你家的水、電、燃氣戶號繫結好。以後每個月收到賬單，直接點進去就能交。戶號一般在你家的紙質賬單上能找到，或者打客服電話問。
  \item \textbf{外賣比價}：同一家店在美團和餓了麼上的價格和優惠經常不一樣。下單前兩邊都看一眼，有時候能省好幾塊。另外，美團那個"拼好飯"功能，和朋友一起點，價格能便宜不少。
  \item \textbf{打車安全}：在滴滴或者高德打車裡，記得把"緊急聯絡人"和"行程分享"設定好。打車的時候，App會自動把你的行程分享給你設定的聯絡人。另外，深夜打車優先選"計程車"而不是"快車"，計程車有行業監管，安全係數更高一些。
\end{itemize}

\booksection{六、用手機搞定醫保、社保、政務}

這幾年的變化很大，以前要跑大廳排隊的很多事，現在手機上幾分鐘就辦完了。

\begin{itemize}
  \item \textbf{醫保電子憑證}：在支付寶或微信裡搜"醫保電子憑證"，啟用之後，去醫院看病、藥店買藥直接刷手機就行，不用帶實體醫保卡。第一次需要人臉識別驗證，之後每次用的時候會生成一個動態碼。
  \item \textbf{社保查詢}：下載"掌上12333"App或者用支付寶裡的"市民中心"，可以查社保繳納記錄、養老金賬戶餘額、失業金申領等。
  \item \textbf{個稅}：每年3月到6月的個稅彙算清繳，在"個人所得稅"App上就能搞定。記得看看有沒有可以申報的專項附加扣除（子女教育、房貸利息、贍養老人等），這些都能幫你少交稅或者退稅。
\end{itemize}

\begin{tcbwarn}{記住這幾條，防騙保平安}
\begin{itemize}
  \item 公檢法、銀行、客服，不會讓你轉賬到"安全賬戶"。
  \item 任何人讓你把付款碼截圖發過去，都是騙子。
  \item 陌生連結別點，尤其是"ETC過期""社保卡停用"這類簡訊。
  \item 網購退款會原路返回，不需要你提供銀行卡號。
\end{itemize}

\end{tcbwarn}

\begin{tcbpractice}{現在就做：給你的支付App做一次安全體檢}
\begin{enumerate}
  \item 開啟支付寶和微信，確認已開啟指紋/面容支付。
  \item 找到"安全鎖"功能並開啟。
  \item 檢查"免密支付/自動扣款"列表，關掉不再需要的服務。
  \item 檢視"登入裝置"列表，移除不認識的裝置。
  \item 啟用你的醫保電子憑證，下次去醫院就不用帶實體卡了。
\end{enumerate}

\end{tcbpractice}
